\subsection{Descriptive results}

Table \ref{tab:top20} reports the twenty firms with the highest AI exposure scores. The group is dominated by semiconductor and software firms for which artificial intelligence was already a core part of the product offering or business strategy by end of 2022. NVIDIA leads with a normalised score of 22.69, more than double the next highest firm, Hewlett Packard Enterprise (10.66). The presence of firms such as Baker Hughes and Ford, which score highly due to AI focused partnerships and subsidiaries rather than internal R\&D, suggests that the measure captures a broad notion of AI engagement that extends beyond the technology sector. The Magnificent Seven firms appear at varying positions: Microsoft (6.82), Alphabet (2.42), Meta (1.78), Amazon (1.73), and Tesla (0.72) all fall within the High AI Exposure group, while Apple and Netflix are classified as having no AI exposure, a point we return to in the limitations discussion below.

\begin{table}[H]
\centering
\caption{Top 20 Firms by AI Exposure Score}
\label{tab:top20}
\small
\begin{tabular}{clcc}
\toprule
Rank & Company & Ticker & AI Score \\
\midrule
1  & NVIDIA             & NVDA  & 22.69 \\
2  & Hewlett Packard Enterprise & HPE & 10.66 \\
3  & Adobe               & ADBE  & 7.91 \\
4  & Intel               & INTC  & 7.79 \\
5  & Microsoft           & MSFT  & 6.82 \\
6  & IBM                 & IBM   & 6.48 \\
7  & Juniper Networks    & JNPR  & 6.34 \\
8  & Baker Hughes        & BKR   & 5.90 \\
9  & Intuit              & INTU  & 3.80 \\
10 & Motorola Solutions  & MSI   & 3.77 \\
11 & Qualcomm            & QCOM  & 3.57 \\
12 & ServiceNow          & NOW   & 3.51 \\
13 & AMD                 & AMD   & 3.03 \\
14 & Autodesk            & ADSK  & 2.96 \\
15 & Salesforce          & CRM   & 2.78 \\
16 & Ford Motor          & F     & 2.77 \\
17 & Zebra Technologies  & ZBRA  & 2.56 \\
18 & Alphabet            & GOOGL & 2.42 \\
19 & Micron Technology   & MU    & 2.36 \\
20 & Cognizant           & CTSH  & 2.28 \\
\bottomrule
\end{tabular}
\vspace{0.4em}
\begin{minipage}{0.95\linewidth}
\footnotesize
\textit{Notes:} AI exposure score is measured as AI-related keyword mentions per 10{,}000 words in the most recent 10-K filing available before 31 December 2022.
\end{minipage}
\end{table}

Table \ref{tab:keywords} reports the frequency of each keyword across all 498 filings. The two most general terms, standalone ``AI'' and ``artificial intelligence,'' account for 934 of the 1{,}278 total mentions (73\%). ``Machine learning'' contributes a further 204 mentions. The more technical terms appear less frequently but help discriminate among the most AI intensive firms. The absence of post-boom terminology such as ``large language model'' or ``generative AI'' from the dictionary is validated by the finding that these terms have near-zero frequency in pre-2023 filings, as confirmed through the expanded keyword screening described above.

\begin{table}[H]
\centering
\caption{Keyword Frequency Across All Pre-2023 Filings}
\label{tab:keywords}
\small
\begin{tabular}{lc}
\toprule
Keyword & Total Mentions \\
\midrule
AI (standalone)              & 611 \\
Artificial intelligence      & 323 \\
Machine learning             & 204 \\
Data science/scientist       & 75 \\
Deep learning                & 23 \\
Neural network(s)            & 18 \\
Computer vision              & 15 \\
Natural language processing  & 9 \\
\midrule
Total                        & 1{,}278 \\
\bottomrule
\end{tabular}
\vspace{0.4em}
\begin{minipage}{0.95\linewidth}
\footnotesize
\textit{Notes:} Counts reflect the total number of keyword appearances across all 498 10-K filings retrieved before 31 December 2022. ``AI (standalone)'' uses word-boundary matching to exclude false positives. ``Neural network(s)'' captures both singular and plural forms. ``Data science/scientist'' captures all variants including ``data sciences,'' ``data scientist,'' and ``data scientists.''
\end{minipage}
\end{table}